\section*{Chapitre \ref{ch:rc:pnl-explain-and-independent-price-verification} --- Explication du P\&L et vérification indépendante des prix}

\begin{solution}{exo:rc:pnl-explain-and-independent-price-verification:1}
$1\,086\times30\times12 = \text{USD}~391\,087$ (avec le vanna non arrondi).
\end{solution}

\begin{solution}{exo:rc:pnl-explain-and-independent-price-verification:2}
Les strikes sont symétriques autour du forward et le modèle est normal~: le payer long et le receiver court
ont un gamma, un vega et un theta égaux et opposés, ainsi qu'un delta égal de même signe.
\end{solution}

\begin{solution}{exo:rc:pnl-explain-and-independent-price-verification:3}
La bande va de 33 à 37~; la marque atteint son bord, 37, soit un ajustement de $-1$.
\end{solution}

\begin{solution}{exo:rc:pnl-explain-and-independent-price-verification:4}
Chaque étape est réévaluée avec les facteurs déjà déplacés, si bien que les effets croisés sont imputés au facteur
qui bouge en dernier. Une moyenne sur tous les ordres (l'allocation de Shapley) supprime la dépendance au prix de
réévaluations supplémentaires~; présenter séparément l'effet croisé constitue l'alternative pratique.
\end{solution}

\begin{solution}{exo:rc:pnl-explain-and-independent-price-verification:5}
Le terme croisé est proportionnel au mouvement des taux~: environ la moitié, soit USD~181 mille pour un mouvement de 15 points de base.
\end{solution}

\begin{solution}{exo:rc:pnl-explain-and-independent-price-verification:6}
Sa juste valeur compense les jambes longue et courte (nulle ici), mais l'incertitude sur le prix de chaque jambe ne
se compense pas~: chacune dépend d'une volatilité dont le consensus est dispersé. L'AVA croît avec le risque brut.
\end{solution}

\begin{solution}{exo:rc:pnl-explain-and-independent-price-verification:7}
Les AVA individuelles doublent (382\,470 et 468\,316) et le montant agrégé atteint USD~425\,393~; le
chiffre simplifié ne change pas.
\end{solution}

\begin{solution}{exo:rc:pnl-explain-and-independent-price-verification:8}
Les seuils autorisent des marques n'importe où à l'intérieur de la bande~; des marques systématiquement placées à son bord favorable passent
tous les tests, comme en 2012. L'IPV ne couvre par ailleurs que les instruments dont les prix sont indépendants~; les positions de niveau~3 et
les intrants de modèle exigent d'autres contrôles, et les réserves ainsi que les AVA restent nécessaires.
\end{solution}

\begin{solution}{pb:rc:pnl-explain-and-independent-price-verification:1}
\textbf{1.} USD~4\,571\,501 réel~; 4\,173\,095 expliqués.
\textbf{2.} USD~398\,406, soit 9~\% du P\&L.
\textbf{3.} USD~139\,103 par point de base~; 1\,086 par point de base de taux et de volatilité.
\textbf{4.} USD~391\,087~; un résiduel de 7\,319.
\textbf{5.} Termes du troisième ordre~: la variation du vanna avec le taux et la volatilité sur un mouvement aussi large.
\textbf{6.} Saisie et données (transactions manquantes, données de marché périmées), puis risques absents (termes croisés,
facteurs non marqués), puis marques (ruptures de l'IPV)~; les deux premières catégories sont les plus rapides à confirmer ou à écarter.
\textbf{7.} Déplacer les facteurs un à un montre que le second facteur déplacé porte un effet qui
dépend du premier~: un effet croisé.
\textbf{8.} L'étape de volatilité ne contribue pas, car le vega est nul au départ~; tout l'effet croisé
apparaît dans l'étape de taux.
\textbf{9.} Le terme croisé est un produit de deux mouvements~; les jours calmes, les deux sont faibles, et leur produit l'est aussi.
\textbf{10.} Une marque plate implique l'absence de skew~; avec le skew du marché, une hausse des taux modifie la volatilité à
chaque strike, et l'explication devrait porter l'effet en vega~; sans cela, il apparaît sous forme de vanna que le desk
ne déclare pas.
\textbf{11.} La volatilité du receiver passe de 95 à 93 points de base~: la valeur du portefeuille augmente de
USD~146\,482.
\textbf{12.} USD~211\,771 selon l'approche de base (190\,578 et 232\,964 agrégés à 50~\%)~; 6\,872 selon l'approche
simplifiée.
\textbf{13.} Un seuil monétaire lié au P\&L quotidien typique du desk et une part (disons quelques pour cent), avec
une investigation automatique les jours de forts mouvements.
\textbf{14.} Non~: la réévaluation s'additionne par construction et masque les risques manquants~; l'\omterm{def:rc:pnl-explain-and-independent-price-verification:risk}{explication par les sensibilités} vérifie
si les mesures de risque du desk décrivent son P\&L.
\textbf{15.} Le \omterm{def:rc:pnl-explain-and-independent-price-verification:unexplained}{P\&L inexpliqué}, sa cause, le vanna absent des rapports de risque du desk, et la correction
(marques de skew et déclaration des grecques croisées).
\textbf{16.} Lorsqu'il ne varie pas avec les mouvements de marché, apparaît les jours calmes ou s'inverse le lendemain~: une
transaction saisie en retard, un notionnel erroné, un fixing périmé.
\textbf{17.} Parce que la rémunération du desk dépend du P\&L qu'il devrait expliquer et des marques qu'il devrait vérifier.
\textbf{18.} Le test d'explication compare le P\&L du modèle de risque à celui des systèmes de valorisation~; un desk dont
le modèle de risque omet le terme croisé présenterait le même écart à cet endroit.
\textbf{19.} Résultat à retenir~: \emph{les 400 mille inexpliqués}~: le terme de vanna manquant, 1\,086 dollars par point de base
au carré multiplié par 30 et 12, explique USD~391\,087 des 398\,406 inexpliqués, laissant 7\,319.
\textbf{20.} Quel risque, quelle donnée ou quelle marque manque à l'image que le desk a de lui-même.
\end{solution}

\begin{solution}{iq:rc:pnl-explain-and-independent-price-verification:1}
Delta, gamma, vega (et volga), theta, les termes croisés pertinents (vanna, \omterm{def:rc:rates-risk:crossgamma}{gamma croisé} entre courbes), carry et financement,
nouvelles transactions et modifications, réserves et frais~; le reste est inexpliqué.
\iqlookfor{grecques, temps, nouvelles affaires et résiduel.}
\end{solution}

\begin{solution}{iq:rc:pnl-explain-and-independent-price-verification:2}
Données de marché manquantes ou périmées, erreurs de saisie, \omterm{def:rc:market-risk-measures:factor}{facteurs de risque} ou termes croisés absents, modifications de modèle,
grands mouvements au-delà de l'approximation, réserves et frais non inclus, et marques ayant changé sans mouvement
de marché.
\iqlookfor{données, saisie, modèle et marques.}
\end{solution}

\begin{solution}{iq:rc:pnl-explain-and-independent-price-verification:3}
Vérifier les marques et les intrants de modèle par rapport à des sources indépendantes et ajuster ceux qui sortent de la tolérance. Cette tâche est
difficile pour les instruments peu liquides et les intrants non observables, lorsque le consensus est mince ou périmé, et lorsque les seuils
incitent à placer les marques à leurs bornes.
\iqlookfor{le processus et ses points faibles.}
\end{solution}

\begin{solution}{iq:rc:pnl-explain-and-independent-price-verification:4}
La juste valeur est un prix de sortie fondé sur les hypothèses des intervenants de marché~; l'\omterm{def:rc:pnl-explain-and-independent-price-verification:pv}{évaluation prudente} retient un point conservateur
(niveau de confiance de 90~\%) de l'intervalle des valeurs de sortie plausibles et déduit l'écart des fonds propres.
\iqlookfor{le niveau de confiance et la déduction en fonds propres.}
\end{solution}

\begin{solution}{iq:rc:pnl-explain-and-independent-price-verification:5}
Grecques d'ouverture par transaction et par facteur, issues de la bibliothèque de valorisation, instantanés de données de marché de fin de journée,
explication par facteur avec seuils et détail, réévaluation de secours pour les grands mouvements, alertes,
piste d'audit et rapprochement avec le grand livre.
\iqlookfor{données, calcul, seuils et contrôles.}
\end{solution}

\begin{solution}{iq:rc:pnl-explain-and-independent-price-verification:6}
Comparer les marques au consensus dans le temps~: la distribution de la position de la marque dans la bande, sa
corrélation avec les objectifs de P\&L du desk, les sauts en fin de mois, et les écarts entre les marques de fin de mois et
les marques intramensuelles.
\iqlookfor{tests statistiques sur la position des marques.}
\end{solution}
